% 51
\begin{frame}{Пример из ML: метод наименьших квадратов}
\small
Рассмотрим уже знакомую задачу
\[
 f(w)=\frac12\norm{Xw-y}^2.
\]
Её гессиан
\[
 \grad^2f(w)=X^\top X.
\]
Для любого $z$
\[
 z^\top X^\top Xz=\norm{Xz}^2\ge0.
\]
Следовательно,
\[
\boxed{X^\top X\succeq0,}
\]
и задача выпукла.
\end{frame}

% 52
\begin{frame}{Когда минимум линейной регрессии единственный?}
\small
Если столбцы $X$ линейно независимы, то
\[
 Xz\ne0\qquad\forall z\ne0.
\]
Поэтому
\[
 z^\top X^\top Xz=\norm{Xz}^2>0,
\]
то есть
\[
 X^\top X\succ0.
\]

Следовательно, функция строго и сильно выпукла, а минимум единственен.
\end{frame}

% 53
\begin{frame}{Ridge добавляет кривизну}
\small
Для Ridge-регрессии
\[
 f(w)=\frac12\norm{Xw-y}^2+\frac{\lambda}{2}\norm{w}^2,
 \qquad \lambda>0,
\]
имеем
\[
 \grad^2f=X^\top X+\lambda I\succeq\lambda I.
\]

\centering
\includegraphics[width=0.70\linewidth]{ridge_geometry.pdf}
\end{frame}

% 54
\begin{frame}{Логистическая потеря тоже выпукла}
\small
Для одного объекта удобно рассматривать
\[
 \ell(z)=\log(1+e^{-z}).
\]
Её вторая производная
\[
 \ell''(z)=\frac{e^z}{(1+e^z)^2}>0.
\]
Следовательно, $\ell$ строго выпукла по $z$.

\centering
\includegraphics[width=0.66\linewidth]{logistic_convexity.pdf}
\end{frame}

% 55
\begin{frame}{Сумма выпуклых функций остаётся выпуклой}
\footnotesize
\begin{block}{Утверждение}
Если $f_i$ выпуклы и $\alpha_i\ge0$, то
\[
 F(x)=\sum_i\alpha_i f_i(x)
\]
выпукла.
\end{block}

Действительно,
\[
\begin{aligned}
F((1-t)x+ty)
&=\sum_i\alpha_i f_i((1-t)x+ty)\\
&\le\sum_i\alpha_i\big((1-t)f_i(x)+tf_i(y)\big)\\
&=(1-t)F(x)+tF(y).
\end{aligned}
\]

Отсюда средний эмпирический риск из выпуклых потерь тоже выпуклый.
\end{frame}

% 56
\begin{frame}{Где эта теория перестаёт работать}
\small
Для параметров глубокой нейронной сети функция потерь обычно невыпукла.

\centering
\includegraphics[width=0.67\linewidth]{nonconvex_loss_landscape.pdf}

\vspace{0.1em}
\small
Поэтому уже нельзя автоматически утверждать:
\[
\text{локальный минимум = глобальный},
\]
\[
\grad f(x)=0\Longrightarrow\text{глобальный минимум}.
\]
\end{frame}

% 57
\begin{frame}{Итоги и следующий вопрос}
\small
Сегодня мы получили цепочку
\[
\boxed{
\text{выпуклость}
\Longrightarrow
f(y)\ge f(x)+\grad f(x)^\top(y-x)
\Longrightarrow
\grad f(x^\star)=0\iff x^\star\text{ — глобальный минимум}
}
\]
и критерий
\[
 f\text{ выпукла}
 \iff
 \grad^2f(x)\succeq0.
\]

Для сильно выпуклой и гладкой функции
\[
 \mu I\preceq\grad^2f(x)\preceq LI,
 \qquad
 \kappa=\frac{L}{\mu}.
\]

\begin{alertblock}{Следующая лекция}
Теперь мы знаем, \textbf{куда} должен прийти градиентный спуск. Следующий вопрос: \textbf{насколько быстро он туда приходит?}
\end{alertblock}
\end{frame}
